%% toc: P2 13 13.4,13.5,13.6,13.7
\section{Proofs for analytic tableau metatheory}
\label{sec:annex-p02-w02-p02-014-tableau-metatheory}

Notation is that of \cref{def:p02-w02-p02-014-input,def:p02-w02-p02-014-expandable}.
For a valuation \(v\) and a set \(X\) of formulas, \(v\entails X\) means
\(v\entails\chi\) for every \(\chi\in X\); for an alternative \(A\),
\(v\entails A\) means the same for the members of \(A\).

\begin{annexitem}{annex:p02-w02-p02-014-alternatives}{thm:p02-w02-p02-014-alternatives}{Each rule is a semantic case split}
\begin{proof}
Fix a valuation \(v\).  Recall from \cref{sec:prop-syntax-semantics-02}
that \(v\entails\neg\alpha\) iff \(\widehat v(\alpha)\) is false.  Treat the
seven shapes in turn.

\emph{\(\varphi\land\psi\).}  The only alternative is
\(\langle\varphi,\psi\rangle\), and \(v\entails\varphi\land\psi\) iff
\(v\entails\varphi\) and \(v\entails\psi\) by the conjunction clause.

\emph{\(\neg(\varphi\lor\psi)\).}  The alternative is
\(\langle\neg\varphi,\neg\psi\rangle\).  Now \(v\entails\neg(\varphi\lor\psi)\)
iff \(\widehat v(\varphi\lor\psi)\) is false iff both \(\widehat v(\varphi)\)
and \(\widehat v(\psi)\) are false iff \(v\entails\neg\varphi\) and
\(v\entails\neg\psi\).

\emph{\(\neg(\varphi\to\psi)\).}  The alternative is
\(\langle\varphi,\neg\psi\rangle\).  The implication clause makes
\(\widehat v(\varphi\to\psi)\) false iff \(\widehat v(\varphi)\) is true and
\(\widehat v(\psi)\) is false, that is, iff \(v\entails\varphi\) and
\(v\entails\neg\psi\).

\emph{\(\neg\neg\varphi\).}  The alternative is \(\langle\varphi\rangle\).
Applying the negation clause twice, \(v\entails\neg\neg\varphi\) iff
\(\widehat v(\neg\varphi)\) is false iff \(\widehat v(\varphi)\) is true.

\emph{\(\neg(\varphi\land\psi)\).}  The alternatives are
\(\langle\neg\varphi\rangle\) and \(\langle\neg\psi\rangle\).
\(\widehat v(\varphi\land\psi)\) is false iff \(\widehat v(\varphi)\) is false
or \(\widehat v(\psi)\) is false, that is, iff \(v\entails\neg\varphi\) or
\(v\entails\neg\psi\).

\emph{\(\varphi\lor\psi\).}  The alternatives are \(\langle\varphi\rangle\)
and \(\langle\psi\rangle\), and the disjunction clause is exactly the
required equivalence.

\emph{\(\varphi\to\psi\).}  The alternatives are \(\langle\neg\varphi\rangle\)
and \(\langle\psi\rangle\).  The implication clause makes
\(\widehat v(\varphi\to\psi)\) true iff \(\widehat v(\varphi)\) is false or
\(\widehat v(\psi)\) is true, that is, iff \(v\entails\neg\varphi\) or
\(v\entails\psi\).

For the size bound, let \(n=\operatorname{sz}\) of the premise and use
\(\operatorname{sz}(\varphi),\operatorname{sz}(\psi)\geq1\).  For
\(\varphi\land\psi\) and \(\varphi\lor\psi\), \(n=\operatorname{sz}(\varphi)+
\operatorname{sz}(\psi)+1\), so \(\varphi\) and \(\psi\) have size at most
\(n-2\).  For \(\neg(\varphi\land\psi)\) and \(\neg(\varphi\lor\psi)\),
\(n=\operatorname{sz}(\varphi)+\operatorname{sz}(\psi)+2\), so
\(\operatorname{sz}(\neg\varphi)=\operatorname{sz}(\varphi)+1\leq n-2\), and
likewise for \(\neg\psi\).  For \(\varphi\to\psi\),
\(\operatorname{sz}(\neg\varphi)=n-\operatorname{sz}(\psi)\leq n-1\) and
\(\operatorname{sz}(\psi)\leq n-2\).  For \(\neg(\varphi\to\psi)\),
\(n=\operatorname{sz}(\varphi)+\operatorname{sz}(\psi)+2\), so
\(\operatorname{sz}(\varphi)\leq n-3\) and
\(\operatorname{sz}(\neg\psi)\leq n-2\).  For \(\neg\neg\varphi\),
\(\operatorname{sz}(\varphi)=n-2\).  In the table of
\cref{def:p02-w02-p02-014-expandable} each alternative has one or two
members.
\end{proof}
\end{annexitem}

\begin{annexitem}{annex:p02-w02-p02-014-invariants}{thm:p02-w02-p02-014-invariants}{Branch invariants}
\begin{proof}
Let \(T_0,\ldots,T_k\) be an expansion sequence ending in the given tableau.
We show by induction on \(j\leq k\) that every branch of \(T_j\) has
properties 1 and 2.

For \(j=0\) the only branch is \(B_\Gamma=(\langle\gamma_1,\ldots,
\gamma_n\rangle,\varnothing)\).  Property 1 holds with \(m=n\), and property 2
is vacuous because no position is marked.

Suppose every branch of \(T_j\) has both properties, and \(T_{j+1}\) replaces
the branch \(B=(\langle\theta_1,\ldots,\theta_m\rangle,M)\) by its children
under an expansion step at position \(i\).  The branches of \(T_{j+1}\) that
are not children of \(B\) are branches of \(T_j\).  A child has the form
\(C=(\langle\theta_1,\ldots,\theta_m\rangle\frown A,\,M\cup\{i\})\) for an
alternative \(A\) of \(\theta_i\).  Its list extends that of \(B\), so its
first \(n\) entries are still \(\gamma_1,\ldots,\gamma_n\), which is
property 1, and \(\operatorname{fm}(B)\subseteq\operatorname{fm}(C)\).  For property 2, let \(i'\) be a marked
position of \(C\).  If \(i'\in M\), the induction hypothesis gives that
\(\theta_{i'}\) is expandable and an alternative \(A'\) of it with
\(A'\subseteq\operatorname{fm}(B)\subseteq\operatorname{fm}(C)\).  If \(i'=i\), then \(\theta_i\) is expandable
because a step is performed only at an expandable formula, and
\(A\subseteq\operatorname{fm}(C)\) by the construction of \(C\).
\end{proof}
\end{annexitem}

\begin{annexitem}{annex:p02-w02-p02-014-soundness}{thm:p02-w02-p02-014-soundness}{Soundness of closure}
\begin{proof}
\emph{Clause 1.}  Let the step be at position \(i\) of \(B\).  Since
\(\theta_i\in\operatorname{fm}(B)\) and \(v\entails\operatorname{fm}(B)\), we have \(v\entails\theta_i\).
By \cref{thm:p02-w02-p02-014-alternatives} some alternative \(A\) of
\(\theta_i\) has \(v\entails A\).  The child \(C\) built from \(A\) has
\(\operatorname{fm}(C)=\operatorname{fm}(B)\cup A\), so \(v\entails\operatorname{fm}(C)\).

\emph{Clause 2.}  Suppose \(v\entails\Gamma\) and let \(T_0,\ldots,T_k\) be
an expansion sequence.  By induction on \(j\), \(T_j\) has a branch \(B\)
with \(v\entails\operatorname{fm}(B)\).  For \(j=0\), \(\operatorname{fm}(B_\Gamma)=\Gamma\).  If \(T_j\) has
such a branch \(B\) and the step from \(T_j\) to \(T_{j+1}\) does not replace
\(B\), then \(B\) is still a branch of \(T_{j+1}\); if it replaces \(B\),
clause 1 gives a child \(C\), a branch of \(T_{j+1}\), with
\(v\entails\operatorname{fm}(C)\).  Any branch \(B\) with \(v\entails\operatorname{fm}(B)\) is open: it does
not contain \(\bottom\), because \(\widehat v(\bottom)\) is false; and it does
not contain both \(\chi\) and \(\neg\chi\), because then
\(\widehat v(\chi)\) would be both true and false.

\emph{Clause 3.}  If \(\Gamma\) were satisfied by some \(v\), clause 2 would
give every tableau rooted at \(\Gamma\) an open branch, so none would be
closed.  Clause 2 does not assume the tableau finished, so clause 3 holds
for all tableaux.
\end{proof}
\end{annexitem}

\begin{annexitem}{annex:p02-w02-p02-014-termination}{thm:p02-w02-p02-014-termination}{Termination of exhaustive expansion}
\begin{proof}
\emph{Clause 1.}  Let the step be at position \(i\) of \(B\), let
\(n=\operatorname{sz}(\theta_i)\), and let \(C\) be the child for the
alternative \(A\).  The positions of \(C\) holding an unmarked expandable
formula are those of \(B\) other than \(i\), which is now marked, together
with those new positions, holding members of \(A\), whose formula is
expandable.  By \cref{thm:p02-w02-p02-014-alternatives}, \(A\) has at most
two members and each has size at most \(n-1\).  Hence
\[
  \mu(C)\leq\mu(B)-3^{n}+2\cdot3^{n-1}=\mu(B)-3^{n-1}<\mu(B).
\]

\emph{Clause 2.}  Record the expansion sequence as a rooted tree: its root
is labelled \(B_\Gamma\), and the step from \(T_j\) to \(T_{j+1}\) adds below
the node of the replaced branch one new node for each child, labelled by that
child.  The branches of each \(T_j\) are exactly the labels of the leaves of
the tree built so far, so every step adds one or two nodes below a leaf.
After \(k\) steps the tree has at least \(k+1\) nodes, and every node has at
most two children.  Along any path from the root the labels have strictly
decreasing \(\mu\) by clause 1.  As \(\mu\) takes values in the natural
numbers and the root has value \(\mu_0\), every node has depth at most
\(\mu_0\).  A rooted tree in which each node has at most two children has at
most \(2^{d}\) nodes at depth \(d\), so this tree has at most
\(\sum_{d=0}^{\mu_0}2^{d}=2^{\mu_0+1}-1\) nodes.  Therefore
\(k+1\leq2^{\mu_0+1}-1\).

\emph{Clause 3.}  If a tableau is not finished, some branch \(B\) is neither
closed nor saturated.  Being open and unsaturated, \(B\) has an unmarked
position holding an expandable formula, and an expansion step is defined
there.  Conversely, if a step is defined at a branch \(B\), then \(B\) is
open and has an unmarked position holding an expandable formula, so \(B\) is
neither closed nor saturated and the tableau is not finished.

\emph{Clause 4.}  Perform steps, in any order, for as long as one is defined.
By clause 2 this stops after at most \(2^{\mu_0+1}-2\) steps, and by clause 3
the tableau at which it stops is finished.  Starting at \(T_0\) produces a
finished tableau rooted at \(\Gamma\).
\end{proof}
\end{annexitem}

\begin{annexitem}{annex:p02-w02-p02-014-truth}{thm:p02-w02-p02-014-truth}{Truth lemma for open saturated branches}
\begin{proof}
Write \(B=(\langle\theta_1,\ldots,\theta_m\rangle,M)\).  We prove, by strong
induction on \(s\geq1\), that \(v_B\entails\chi\) for every \(\chi\in\operatorname{fm}(B)\)
with \(\operatorname{sz}(\chi)=s\).  Fix such a \(\chi\) and assume the claim
for all members of \(\operatorname{fm}(B)\) of smaller size.  By the classification in
\cref{def:p02-w02-p02-014-expandable}, \(\chi\) is a variable, \(\bottom\),
a negated variable, \(\neg\bottom\), or expandable.

\emph{\(\chi=p\).}  Then \(p\in\operatorname{fm}(B)\), so \(v_B(p)\) is true by the
definition of \(v_B\).

\emph{\(\chi=\bottom\).}  This case does not occur, because \(B\) is open.

\emph{\(\chi=\neg p\).}  Since \(B\) is open and \(\neg p\in\operatorname{fm}(B)\), the
variable \(p\) is not in \(\operatorname{fm}(B)\).  So \(v_B(p)\) is false and
\(v_B\entails\neg p\).

\emph{\(\chi=\neg\bottom\).}  \(\widehat{v_B}(\bottom)\) is false, so
\(v_B\entails\neg\bottom\).

\emph{\(\chi\) expandable.}  Let \(i\) be a position with
\(\theta_i=\chi\).  As \(B\) is saturated, \(i\in M\).  By
\cref{thm:p02-w02-p02-014-invariants}(2) there is an alternative \(A\) of
\(\chi\) with \(A\subseteq\operatorname{fm}(B)\).  By
\cref{thm:p02-w02-p02-014-alternatives} each member of \(A\) has size less
than \(s\), so the induction hypothesis gives \(v_B\entails A\).  The
right-to-left direction of \cref{thm:p02-w02-p02-014-alternatives} then gives
\(v_B\entails\chi\).

These cases exhaust the formulas, which completes the induction.
\end{proof}
\end{annexitem}

\begin{annexitem}{annex:p02-w02-p02-014-completeness}{thm:p02-w02-p02-014-completeness}{Completeness and read-off}
\begin{proof}
\emph{Clause 2.}  Let \(B\) be an open branch of the finished tableau \(T\).
Since \(T\) is finished and \(B\) is not closed, \(B\) is saturated.  By
\cref{thm:p02-w02-p02-014-truth}, \(v_B\entails\operatorname{fm}(B)\), and by
\cref{thm:p02-w02-p02-014-invariants}(1), \(\Gamma\subseteq\operatorname{fm}(B)\), so
\(v_B\entails\Gamma\).

\emph{Clause 1.}  If \(T\) is closed, \(\Gamma\) is unsatisfiable by
\cref{thm:p02-w02-p02-014-soundness}(3).  If \(T\) is not closed, it has an
open branch \(B\), and clause 2 shows that \(v_B\) satisfies \(\Gamma\).

\emph{Clause 3.}  By clause 1 each finished tableau rooted at \(\Gamma\) is
closed exactly when \(\Gamma\) is unsatisfiable, a condition that does not
mention the tableau.
\end{proof}
\end{annexitem}

\begin{annexitem}{annex:p02-w02-p02-014-consequence}{thm:p02-w02-p02-014-consequence}{Consequence and validity tests}
\begin{proof}
First, \(\Gamma\entails\varphi\) iff \(\Gamma\cup\{\neg\varphi\}\) is
unsatisfiable.  If \(v\) satisfied \(\Gamma\cup\{\neg\varphi\}\) while
\(\Gamma\entails\varphi\), then \(v\entails\Gamma\) would give
\(v\entails\varphi\), contradicting \(v\entails\neg\varphi\).  Conversely, if
\(\Gamma\entails\varphi\) fails, some \(v\) satisfies \(\Gamma\) with
\(\widehat v(\varphi)\) false; then \(v\entails\neg\varphi\), so \(v\)
satisfies \(\Gamma\cup\{\neg\varphi\}\).

\emph{Clause 1.}  By \cref{thm:p02-w02-p02-014-termination}(4) a finished
tableau rooted at \(\Gamma\cup\{\neg\varphi\}\) exists.  By
\cref{thm:p02-w02-p02-014-completeness}(1) any one of them is closed iff
\(\Gamma\cup\{\neg\varphi\}\) is unsatisfiable, iff
\(\Gamma\entails\varphi\).  So ``some'' and ``every'' agree.

\emph{Clause 2.}  By \cref{thm:p02-w02-p02-014-completeness}(2),
\(v_B\entails\Gamma\cup\{\neg\varphi\}\).  Thus \(v_B\entails\Gamma\) and
\(\widehat{v_B}(\varphi)\) is false.

\emph{Clause 3.}  Every valuation satisfies \(\varnothing\), so
\(\varnothing\entails\varphi\) iff \(\entails\varphi\), and
\(\varnothing\cup\{\neg\varphi\}=\{\neg\varphi\}\).  Apply clause 1 with
\(\Gamma=\varnothing\).
\end{proof}
\end{annexitem}

\begin{annexitem}{annex:p02-w02-p02-014-decision}{thm:p02-w02-p02-014-decision}{Propositional tableaux decide by terminating saturation}
\begin{proof}
Each iteration of step 2 of the procedure is applied to a branch that is open
and not saturated, at an unmarked position holding an expandable formula, so
it is an expansion step and the run of the procedure is an expansion sequence
from \(B_\Gamma\).  By \cref{thm:p02-w02-p02-014-termination}(2) it has at
most \(2^{\mu(B_\Gamma)+1}-2\) steps.  The loop exits only when no branch is
open and unsaturated, that is, at a finished tableau \(T\).

If every branch of \(T\) is closed, \(\Gamma\) is unsatisfiable by
\cref{thm:p02-w02-p02-014-soundness}(3).  If \(\Gamma\) is unsatisfiable,
\(T\) is closed by \cref{thm:p02-w02-p02-014-completeness}(1), so the output
is \emph{unsatisfiable}.  Otherwise the output is the set of variables in
\(\operatorname{fm}(B)\) for an open branch \(B\) of \(T\); this finite set determines
\(v_B\), and \(v_B\entails\Gamma\) by
\cref{thm:p02-w02-p02-014-completeness}(2).

Every operation of the procedure acts on finite lists of finite strings:
recognising the shape of a formula, which is unique by
\cref{thm:prop-unique-readability}; forming the alternatives; testing
closure by comparing the finitely many formulas on a branch with \(\bottom\)
and with one another's negations; and testing saturation by inspecting the
finitely many positions.  So the procedure is an algorithm that halts on
every finite \(\Gamma\) with the correct answer, which decides
satisfiability of finite sets.  Running it on \(\Gamma\cup\{\neg\varphi\}\)
and on \(\{\neg\varphi\}\) decides \(\Gamma\entails\varphi\) and
\(\entails\varphi\) by \cref{thm:p02-w02-p02-014-consequence}(1) and (3), and
when the answer is negative the output valuation is a countervaluation by
\cref{thm:p02-w02-p02-014-consequence}(2).
\end{proof}
\end{annexitem}

\begin{annexitem}{annex:p02-w02-p02-014-fol-no-finite-saturation}{thm:p02-w02-p02-014-fol-no-finite-saturation}{No finite first-order saturation}
\begin{proof}
Let the occurrence \(o\) of \(\forall x\,\varphi\) (or of
\(\neg\exists x\,\varphi\)) lie on the finite open branch \(B\).  By
\cref{def:p02-fol-tableau} and the \(\forall\) and \(\neg\exists\) rows of the calculus
sheet \texttt{sheet:\allowbreak tableaux-fol}, a closed
\(\mathcal L_{\Gamma}^{*}\)-term \(t\) is used for \(o\) on \(B\) when the
gamma expansion of \(o\) with \(t\) has been performed there, and each such
expansion appends its instance \(\varphi[t/x]\) (or
\(\neg(\varphi[t/x])\)) at a new position of \(B\).  Distinct uses occupy
distinct positions, even when the instances coincide as formulas.  As \(B\)
has finitely many positions, only finitely many terms are used for \(o\) on
\(B\).

The signature \(\mathcal L_{\Gamma}^{*}\) contains the pairwise distinct
parameter constants \(a_n\) for every \(n\in\mathbb N\), each a closed term,
so it has infinitely many closed terms, and every one of them occurs in the
fixed enumeration.  Choose one, \(t\), not used for \(o\) on \(B\).  Since
\(B\) is open, the gamma expansion of \(o\) with \(t\) is applicable on \(B\),
and it has not been performed.  By \cref{def:p02-fol-tableau}, a branch is
saturated only if every expansion applicable on it has been made, so \(B\) is not
saturated.  The second sentence of the lemma is the contrapositive of the
first.
\end{proof}
\end{annexitem}

\subsection{First-order tableaux}
\label{sec:annex-p02-w02-p02-014-fol}

Notation is that of
\cref{def:p02-w02-p02-014-fol-setting,def:p02-w02-p02-014-fol-equality,def:p02-w02-p02-014-fol-branch,def:p02-w02-p02-014-hintikka}.
For a sentence \(\sigma\), \(\mathcal M\entails\sigma\) means
\(\mathcal M,s\entails\sigma\) for one, equivalently every, assignment \(s\);
for a set \(X\) of sentences, \(\mathcal M\entails X\) means
\(\mathcal M\entails\sigma\) for every \(\sigma\in X\).  For a closed term
\(t\), \(t^{\mathcal M}=\sem{t}^{\mathcal M}_s\) for every \(s\), by term
coincidence.  Two facts are used repeatedly.  First, by
\cref{thm:fol-substitution}, for every formula \(\varphi\) whose free
variables are among \(x\), every closed term \(t\) and every assignment \(s\),
\[
  \mathcal M\entails\varphi[t/x]
  \quad\text{iff}\quad
  \mathcal M,s[x\mapsto t^{\mathcal M}]\entails\varphi .
  \tag{S}
\]
Second, by \cref{lem:reduct-invariance}, if \(\mathcal L''\) extends
\(\mathcal L'\), the truth of an \(\mathcal L'\)-formula under an assignment
in an \(\mathcal L''\)-structure depends only on its reduct to
\(\mathcal L'\).

\begin{annexitem}{annex:p02-w02-p02-014-fol-cases}{thm:p02-w02-p02-014-fol-cases}{Sound case splits}
\begin{proof}
\emph{Clause 1.}  For sentences the satisfaction clauses give:
\(\mathcal M\not\entails\bottom\); \(\mathcal M\entails\neg\alpha\) iff
\(\mathcal M\not\entails\alpha\); \(\mathcal M\entails\alpha\land\beta\) iff
both conjuncts hold; \(\mathcal M\entails\alpha\lor\beta\) iff at least one
disjunct holds; \(\mathcal M\entails\alpha\to\beta\) iff
\(\mathcal M\not\entails\alpha\) or \(\mathcal M\entails\beta\).  Hence:
\(\alpha\land\beta\) holds iff \(\alpha\) and \(\beta\) hold;
\(\neg(\alpha\lor\beta)\) holds iff neither disjunct holds, iff
\(\neg\alpha\) and \(\neg\beta\) hold; \(\neg(\alpha\to\beta)\) holds iff
\(\alpha\) holds and \(\beta\) fails, iff \(\alpha\) and \(\neg\beta\) hold;
\(\neg\neg\alpha\) holds iff \(\neg\alpha\) fails, iff \(\alpha\) holds;
\(\neg(\alpha\land\beta)\) holds iff a conjunct fails, iff \(\neg\alpha\) or
\(\neg\beta\) holds; \(\alpha\lor\beta\) holds iff \(\alpha\) or \(\beta\)
holds; and \(\alpha\to\beta\) holds iff \(\neg\alpha\) or \(\beta\) holds.
These are the seven rows of the table of alternatives in
\cref{def:p02-w02-p02-014-expandable}, read for \(\mathcal M\).

\emph{Clause 2.}  If \(\mathcal M\entails\forall x\,\varphi\), then
\(\mathcal M,s[x\mapsto d]\entails\varphi\) for every element \(d\), in
particular for \(d=t^{\mathcal M}\), so \(\mathcal M\entails\varphi[t/x]\) by
(S).  If \(\mathcal M\entails\neg\exists x\,\varphi\), no element \(d\) has
\(\mathcal M,s[x\mapsto d]\entails\varphi\); taking \(d=t^{\mathcal M}\) and
using (S), \(\mathcal M\not\entails\varphi[t/x]\), that is
\(\mathcal M\entails\neg(\varphi[t/x])\).

\emph{Clause 3.}  Let \(\mathcal L'\) be \(\mathcal L_{\Gamma}^{*}\) without
\(a\).  Every member of \(X\) is an \(\mathcal L'\)-sentence, and so is
\(\exists x\,\varphi\); hence \(\varphi\) is an \(\mathcal L'\)-formula.  As
\(\mathcal M\entails\exists x\,\varphi\), choose an element \(d\) with
\(\mathcal M,s[x\mapsto d]\entails\varphi\).  Let \(\mathcal M'\) have the
domain and interpretations of \(\mathcal M\) except that
\(a^{\mathcal M'}=d\).  Then \(\mathcal M'\) and \(\mathcal M\) have the same
reduct to \(\mathcal L'\), so by reduct invariance \(\mathcal M'\entails X\)
and \(\mathcal M',s[x\mapsto d]\entails\varphi\).  Since
\(a^{\mathcal M'}=d\), (S) gives \(\mathcal M'\entails\varphi[a/x]\).  For
\(\neg\forall x\,\varphi\in X\), choose \(d\) with
\(\mathcal M,s[x\mapsto d]\not\entails\varphi\); the same \(\mathcal M'\)
satisfies \(X\) and has \(\mathcal M',s[x\mapsto d]\not\entails\varphi\), so
\(\mathcal M'\not\entails\varphi[a/x]\) by (S), that is
\(\mathcal M'\entails\neg(\varphi[a/x])\).

\emph{Clause 4.}  \(\mathcal M\entails t=t\) because
\(t^{\mathcal M}=t^{\mathcal M}\).  If \(\mathcal M\entails r=u\), then
\(r^{\mathcal M}=u^{\mathcal M}\).  The only free variable of \(\varphi\) is
\(x_0\), so by (S) applied twice,
\(\mathcal M\entails\varphi[r/x_0]\) iff
\(\mathcal M,s[x_0\mapsto r^{\mathcal M}]\entails\varphi\) iff
\(\mathcal M,s[x_0\mapsto u^{\mathcal M}]\entails\varphi\) iff
\(\mathcal M\entails\varphi[u/x_0]\).
\end{proof}
\end{annexitem}

\begin{annexitem}{annex:p02-w02-p02-014-fol-invariants}{thm:p02-w02-p02-014-fol-invariants}{First-order branch invariants}
\begin{proof}
\emph{Finite branches.}  Induct along an expansion sequence
\(T_0,T_1,\ldots\).  The root \(B_\Gamma\) has \(\operatorname{fm}(B_\Gamma)=\Gamma\) and
no done task.  Suppose every branch of \(T_j\) has the stated properties and
\(T_{j+1}\) replaces \(B=(\langle\theta_1,\ldots,\theta_m\rangle,D)\) at the
task \(\tau\) by children
\(C=(\langle\theta_1,\ldots,\theta_m\rangle\frown A,D\cup\{\tau\})\).  The
other branches are unchanged.  For each child,
\(\Gamma\subseteq\operatorname{fm}(B)\subseteq\operatorname{fm}(C)\) and the positions of \(B\) keep their
sentences, so every task in \(D\) still has its conclusion on \(C\).  The
conclusion of \(\tau\) is in \(A\subseteq\operatorname{fm}(C)\): by
\cref{def:p02-w02-p02-014-fol-branch}(3), \(A\) is an alternative of
\(\theta_i\) for a propositional \(\langle i\rangle\), the instance
\(\varphi[a/x]\) or \(\neg(\varphi[a/x])\) at the chosen parameter \(a\) for a
delta \(\langle i\rangle\), the instance of \(\theta_i\) at \(t\) for
\(\langle i,t\rangle\), \(t=t\) for \(\langle{=},t\rangle\), and
\(\varphi[u/x_0]\) for \(\langle i,j,\varphi\rangle\).

In \(\mathsf T\), \(\Gamma\) has no equation and only the sheet rules apply.
An appended sentence is a subformula of the premise, the negation of one,
or an instance \(\varphi[t/x]\) or \(\neg(\varphi[t/x])\) of the
matrix \(\varphi\) of a quantified premise or of a quantified sentence under
a negation.  Substitution of a term replaces terms inside atomic formulas,
keeping their relation symbol, and commutes with or renames through every
other constructor (\cref{sec:fol-syntax-semantics-04}), so it introduces no
equality sign.  By induction along the sequence no sentence on any branch
contains one.

\emph{Limit branches.}  Let \((B_j)\) be a path with limit \(B_\omega\).  Each
list extends the previous one, so a position of \(B_\omega\) holds the same
sentence on every \(B_j\) long enough to contain it.  We have
\(\Gamma\subseteq\operatorname{fm}(B_0)\subseteq\operatorname{fm}(B_\omega)\).  A done task of \(B_\omega\)
lies in the done set of some \(B_j\), so its conclusion is in
\(\operatorname{fm}(B_j)\subseteq\operatorname{fm}(B_\omega)\); in \(\mathsf T\), every member of
\(\operatorname{fm}(B_\omega)\) lies in some \(\operatorname{fm}(B_j)\) and has no equality sign.  If
\(B_\omega\) is closed, the witness \(\bottom\), or the pair \(\chi,\neg\chi\),
occupies at most two positions, which all lie on some \(B_j\), so \(B_j\) is
closed.  If some \(B_j\) is closed, then \(\operatorname{fm}(B_j)\subseteq\operatorname{fm}(B_\omega)\) makes
\(B_\omega\) closed.
\end{proof}
\end{annexitem}

\begin{annexitem}{annex:p02-w02-p02-014-fol-soundness}{thm:p02-w02-p02-014-fol-soundness}{First-order soundness of closure}
\begin{proof}
\emph{Clause 1.}  Let \(\mathcal M\entails\operatorname{fm}(B)\) and let the step be at
\(\tau\).  If \(\tau=\langle i\rangle\) with \(\theta_i\)
propositional-expandable, \cref{thm:p02-w02-p02-014-fol-cases}(1) gives an
alternative \(A\) with \(\mathcal M\entails A\), and the child \(C\) for \(A\)
has \(\operatorname{fm}(C)=\operatorname{fm}(B)\cup A\).  If \(\theta_i\) is a delta sentence, the step
uses a parameter \(a\) occurring in no \(\theta_k\), and
\cref{thm:p02-w02-p02-014-fol-cases}(3) with \(X=\operatorname{fm}(B)\) gives
\(\mathcal M'\) satisfying \(\operatorname{fm}(B)\) and the appended instance.  For
\(\langle i,t\rangle\) use \cref{thm:p02-w02-p02-014-fol-cases}(2).  For
\(\langle{=},t\rangle\) and \(\langle i,j,\varphi\rangle\) use
\cref{thm:p02-w02-p02-014-fol-cases}(4), with \(r=u\) and
\(\varphi[r/x_0]=\theta_j\) both in \(\operatorname{fm}(B)\).

\emph{Clause 2.}  Let \(\mathcal N\) be an \(\mathcal L\)-structure with
\(\mathcal N\entails\Gamma\).  The signature \(\mathcal L\) extends
\(\mathcal L_{\Gamma}\), so the reduct \(\mathcal N\restriction\mathcal
L_{\Gamma}\) satisfies \(\Gamma\) by reduct invariance.  Fix an element
\(e\) of its nonempty domain and expand it to an
\(\mathcal L_{\Gamma}^{*}\)-structure \(\mathcal M\) with
\(a_n^{\mathcal M}=e\) for every \(n\).  Reduct invariance again gives
\(\mathcal M\entails\Gamma=\operatorname{fm}(B_\Gamma)\).  By induction along an expansion
sequence, using clause 1 at the replaced branch and keeping the same
structure for the unchanged branches, every \(T_j\) has a branch \(B\) and an
\(\mathcal L_{\Gamma}^{*}\)-structure satisfying \(\operatorname{fm}(B)\).  Such a branch
is open: \(\bottom\) is true in no structure, and \(\chi\) and \(\neg\chi\)
are not both true in one.

\emph{Clause 3.}  If \(\Gamma\) had a model, clause 2 would give every finite
tableau rooted at \(\Gamma\) an open branch.
\end{proof}
\end{annexitem}

\begin{annexitem}{annex:p02-w02-p02-014-congruence}{thm:p02-w02-p02-014-congruence}{The branch equality is a congruence}
\begin{proof}
If \(r\sim_H u\) with \(r\) and \(u\) different terms, then \(r=u\in H\), so
\(H\) contains the equality sign and the two closure conditions of (H4) are
available.  Every formula \(\varphi\) used below is atomic, contains the
variable \(x_0\) exactly once and contains no other variable, because all
other terms in it are closed; so it is admissible in (H4).

\emph{Equivalence.}  Reflexivity holds by the identity clause.  For
symmetry, let \(r\sim_H u\) with \(r,u\) different, so \(r=u\in H\).  By (H4),
\(r=r\in H\).  With \(\varphi\) the formula \(x_0=r\), we have
\(\varphi[r/x_0]=(r=r)\in H\), so (H4) gives \(\varphi[u/x_0]=(u=r)\in H\)
and \(u\sim_H r\).  For transitivity, let \(r\sim_H u\) and \(u\sim_H w\).
If either pair consists of one term twice, the claim is immediate.
Otherwise \(r=u\in H\) and \(u=w\in H\); with \(\varphi\) the formula
\(r=x_0\), \(\varphi[u/x_0]=(r=u)\in H\), so (H4), applied with the equation
\(u=w\), gives \(\varphi[w/x_0]=(r=w)\in H\) and \(r\sim_H w\).

\emph{Functions.}  Let \(t_i\sim_H u_i\) for \(i\leq n\), and for
\(0\leq i\leq n\) let \(g_i=f(u_1,\ldots,u_i,t_{i+1},\ldots,t_n)\).  We show
\(g_{i-1}\sim_H g_i\) for \(1\leq i\leq n\); transitivity then gives
\(g_0\sim_H g_n\).  If \(t_i\) and \(u_i\) are the same term, so are
\(g_{i-1}\) and \(g_i\).  Otherwise \(t_i=u_i\in H\), and (H4) gives
\(g_{i-1}=g_{i-1}\in H\).  Let \(\varphi\) be the formula
\[
  g_{i-1}=f(u_1,\ldots,u_{i-1},x_0,t_{i+1},\ldots,t_n).
\]
Then
\(\varphi[t_i/x_0]\) is \(g_{i-1}=g_{i-1}\), in \(H\), so (H4) gives
\(\varphi[u_i/x_0]\), which is \(g_{i-1}=g_i\), in \(H\).

\emph{Relations.}  Let \(t_i\sim_H u_i\) and \(R(t_1,\ldots,t_n)\in H\), and
let
\[
  \rho_i=R(u_1,\ldots,u_i,t_{i+1},\ldots,t_n)\qquad(0\leq i\leq n).
\]
By induction on \(i\),
\(\rho_i\in H\): \(\rho_0\in H\) by assumption; if \(t_i\) and \(u_i\) are
the same, \(\rho_i=\rho_{i-1}\); otherwise \(t_i=u_i\in H\), and with
\(\varphi\) the formula
\[
  R(u_1,\ldots,u_{i-1},x_0,t_{i+1},\ldots,t_n),
\]
\(\varphi[t_i/x_0]=\rho_{i-1}\in H\) gives
\(\varphi[u_i/x_0]=\rho_i\in H\).  So \(\rho_n=R(u_1,\ldots,u_n)\in H\).

\emph{Consequences.}  The domain is nonempty because \(a_0\in\mathrm{CT}_{\Gamma}\).
By the function clause, \([f(t_1,\ldots,t_n)]\) depends only on the classes
\([t_i]\), so \(f^{\mathcal M_H}\) is well defined; \(R^{\mathcal M_H}\) is a
set of tuples of classes and needs no choice.  By induction on the closed
term \(t=f(t_1,\ldots,t_n)\), with \(n=0\) for constants including the
parameters, \(t^{\mathcal M_H}=f^{\mathcal M_H}([t_1],\ldots,[t_n])=[t]\).  If
\(([t_1],\ldots,[t_n])\in R^{\mathcal M_H}\), the tuple equals
\(([t'_1],\ldots,[t'_n])\) for some \(R(t'_1,\ldots,t'_n)\in H\); since
\(t'_i\sim_H t_i\), the relation clause gives \(R(t_1,\ldots,t_n)\in H\).  The
converse is the definition of \(R^{\mathcal M_H}\).

If \(H\) is equality-free, no equation is in \(H\), so \(r\sim_H u\) only when
\(r\) and \(u\) are the same term.  Then \([t]=\{t\}\); identifying \([t]\)
with \(t\), the domain is \(\mathrm{CT}_{\Gamma}\), each \(f\) acts by term
formation, and \(R\) holds of \((t_1,\ldots,t_n)\) exactly when
\(R(t_1,\ldots,t_n)\in H\), which is the description of
\(\mathcal M_{B^{+}_{\mathrm{at}}}\) in \cref{thm:p02-fol-tableau-adequacy}.
\end{proof}
\end{annexitem}

\begin{annexitem}{annex:p02-w02-p02-014-herbrand-truth}{thm:p02-w02-p02-014-herbrand-truth}{Herbrand truth lemma}
\begin{proof}
Write \(\mathcal M=\mathcal M_H\).  For a sentence \(\chi\) let
\(\rho(\chi)=2c(\chi)\) if \(\chi\) is a negation and
\(\rho(\chi)=2c(\chi)+1\) otherwise.  We prove \(\mathcal M\entails\chi\) for
every \(\chi\in H\) by strong induction on \(\rho(\chi)\).  First, every
sentence that (H1), (H2) or (H3) requires to be in \(H\) because
\(\chi\in H\) has smaller \(\rho\) than \(\chi\).  If \(\chi\) is a negation,
each such sentence \(\beta\) has \(c(\beta)\leq c(\chi)-1\): for
\(\neg(\alpha\land\beta')\) and \(\neg(\alpha\lor\beta')\) the sentences
\(\neg\alpha,\neg\beta'\) have complexity at most
\(c(\alpha)+c(\beta')+1\); for \(\neg(\alpha\to\beta')\), \(\alpha\) and
\(\neg\beta'\) have complexity at most \(c(\alpha)+c(\beta')+1\); for
\(\neg\neg\alpha\), \(c(\alpha)=c(\chi)-2\); and for
\(\neg\forall x\,\varphi\) and \(\neg\exists x\,\varphi\), the instance
\(\neg(\varphi[t/x])\) has complexity \(c(\varphi)+1=c(\chi)-1\) by
\cref{lem:substitution-complexity}.  Hence
\(\rho(\beta)\leq2c(\chi)-1<\rho(\chi)\).  If \(\chi\) is not a negation, the
required sentences are immediate subformulas or instances
\(\varphi[t/x]\), of complexity at most \(c(\chi)-1\) and so of
\(\rho\leq2c(\chi)-1\); the one exception is \(\neg\alpha\) for
\(\chi=\alpha\to\beta'\), where \(\rho(\neg\alpha)=2c(\alpha)+2\leq
2c(\chi)<\rho(\chi)\).

Fix \(\chi\in H\) and assume the claim for members of \(H\) of smaller
\(\rho\).  By \cref{def:p02-w02-p02-014-fol-setting}, \(\chi\) is of exactly
one kind.

\emph{Atomic.}  If \(\chi=R(t_1,\ldots,t_n)\), then
\(([t_1],\ldots,[t_n])\in R^{\mathcal M}\) and
\(t_i^{\mathcal M}=[t_i]\) by \cref{thm:p02-w02-p02-014-congruence}, so
\(\mathcal M\entails\chi\).  If \(\chi\) is \(r=u\), then \(r\sim_H u\), so
\(r^{\mathcal M}=[r]=[u]=u^{\mathcal M}\).  The case \(\chi=\bottom\) does not
occur by (H0).

\emph{Negated atomic, and \(\neg\bottom\).}  \(\mathcal M\entails\neg\bottom\)
holds.  If \(\chi=\neg R(t_1,\ldots,t_n)\) and
\(\mathcal M\entails R(t_1,\ldots,t_n)\), then
\(([t_1],\ldots,[t_n])\in R^{\mathcal M}\), so \(R(t_1,\ldots,t_n)\in H\) by
\cref{thm:p02-w02-p02-014-congruence}, against (H0).  If \(\chi\) is
\(\neg(r=u)\) and \([r]=[u]\), then \(r\sim_H u\).  Either \(r=u\in H\),
against (H0), or \(r\) and \(u\) are the same term; then \(\chi\) contains the
equality sign, (H4) gives \(r=r\in H\), and \(\chi\) is \(\neg(r=r)\),
against (H0).  So \(\mathcal M\entails\chi\) in each case.

\emph{Propositional-expandable.}  By (H1) every member of some alternative
\(A\) of \(\chi\) is in \(H\), with smaller \(\rho\); by the induction
hypothesis \(\mathcal M\) satisfies them, and
\cref{thm:p02-w02-p02-014-fol-cases}(1) gives \(\mathcal M\entails\chi\).

\emph{Gamma.}  Let \(\chi=\forall x\,\varphi\).  By (H2) and the induction
hypothesis, \(\mathcal M\entails\varphi[t/x]\) for every
\(t\in\mathrm{CT}_{\Gamma}\), so by (S)
\(\mathcal M,s[x\mapsto[t]]\entails\varphi\).  Every element of the domain
is some \([t]\), so \(\mathcal M\entails\forall x\,\varphi\).  Let
\(\chi=\neg\exists x\,\varphi\).  By (H2) and the induction hypothesis,
\(\mathcal M\not\entails\varphi[t/x]\), so
\(\mathcal M,s[x\mapsto[t]]\not\entails\varphi\), for every \(t\); hence
\(\mathcal M\not\entails\exists x\,\varphi\) and \(\mathcal M\entails\chi\).

\emph{Delta.}  Let \(\chi=\exists x\,\varphi\).  By (H3) some
\(\varphi[t/x]\in H\); by the induction hypothesis and (S),
\(\mathcal M,s[x\mapsto[t]]\entails\varphi\), so \(\mathcal M\entails\chi\).
Let \(\chi=\neg\forall x\,\varphi\).  By (H3) some \(\neg(\varphi[t/x])\in H\);
by the induction hypothesis \(\mathcal M\not\entails\varphi[t/x]\), so
\(\mathcal M,s[x\mapsto[t]]\not\entails\varphi\) by (S), and
\(\mathcal M\entails\chi\).
\end{proof}
\end{annexitem}

\begin{annexitem}{annex:p02-w02-p02-014-saturated-hintikka}{thm:p02-w02-p02-014-saturated-hintikka}{Open saturated branches are Hintikka sets}
\begin{proof}
Let \(H=\operatorname{fm}(B)\).  (H0) holds because \(B\) is open.  Let \(\chi\in H\) sit
at position \(i\).  If \(\chi\) is propositional-expandable or delta, the task
\(\langle i\rangle\) is available, hence done by saturation, and
\cref{thm:p02-w02-p02-014-fol-invariants} puts an alternative of \(\chi\),
or an instance at a parameter, on \(B\); this is (H1) or (H3), since every
parameter is in \(\mathrm{CT}_{\Gamma}\).  If \(\chi\) is a gamma sentence,
every task \(\langle i,t\rangle\) is available and done, which puts the
instance at every \(t\) on \(B\); this is (H2).  In \(\mathsf T\),
\cref{thm:p02-w02-p02-014-fol-invariants} shows that \(H\) is equality-free,
so (H4) holds vacuously.  In \(\mathsf T^{=}\), every task
\(\langle{=},t\rangle\) is available and done, so \(t=t\in H\); and if
\(r=u\) is at position \(i\) and \(\varphi[r/x_0]\) is at position \(j\), the
task \(\langle i,j,\varphi\rangle\) is available and done, so
\(\varphi[u/x_0]\in H\).  This is (H4).
\end{proof}
\end{annexitem}

\begin{annexitem}{annex:p02-w02-p02-014-fol-extraction}{thm:p02-w02-p02-014-fol-extraction}{Extraction}
\begin{proof}
By \cref{thm:p02-w02-p02-014-saturated-hintikka}, \(H\) is a Hintikka set,
so \cref{thm:p02-w02-p02-014-herbrand-truth} gives \(\mathcal M_H\entails\chi\)
for every \(\chi\in H\).  By \cref{thm:p02-w02-p02-014-fol-invariants},
\(\Gamma\subseteq H\).  The signature \(\mathcal L_{\Gamma}^{*}\) extends
\(\mathcal L_{\Gamma}\), so reduct invariance gives
\(\mathcal M_H\restriction\mathcal L_{\Gamma}\entails\Gamma\).  The signature
\(\mathcal L\) extends \(\mathcal L_{\Gamma}\); give each symbol of
\(\mathcal L\) outside \(\mathcal L_{\Gamma}\) any interpretation of its
arity on the nonempty domain.  The resulting \(\mathcal L\)-structure has
reduct \(\mathcal M_H\restriction\mathcal L_{\Gamma}\) to \(\mathcal
L_{\Gamma}\), so it satisfies \(\Gamma\) by reduct invariance.  In
\(\mathsf T\), \(H\) is equality-free, and the last clause of
\cref{thm:p02-w02-p02-014-congruence} identifies \(\mathcal M_H\) with
\(\mathcal M_{B^{+}_{\mathrm{at}}}\).
\end{proof}
\end{annexitem}

\begin{annexitem}{annex:p02-w02-p02-014-fol-semidecision}{thm:p02-w02-p02-014-fol-semidecision}{Fair systematic expansion semidecides unsatisfiability}
\begin{proof}
\emph{Clause 1.}  Fix stage \(n\) and an open branch \(B\) of \(T^{(n)}\) of
length \(m_B\).  The list \(W_n(B)\) is finite: it has at most \(m_B\) tasks
\(\langle i\rangle\), at most \(m_B(n+1)\) tasks \(\langle i,t_k\rangle\),
\(n+1\) reflexivity tasks, and, for each of the at most \(m_B^2\) pairs of
positions, one replacement task for each occurrence of \(r\) as a subterm of
\(\theta_j\).  Each task of \(W_n(B)\) is available on every descendant of
\(B\), because descendants extend the list of \(B\) and keep its sentences at
positions up to \(m_B\).  Processing one task \(\tau\) is a run of expansion
steps, each legal because it is made at an open branch whose done set lacks
\(\tau\).  Each such step replaces one open descendant of \(B\) lacking
\(\tau\) by children that have \(\tau\) done, so the finite number of open
descendants lacking \(\tau\) decreases at every step, and processing \(\tau\)
ends.  Branches that descend from different branches of \(T^{(n)}\) are not
affected by each other's steps.  So stage \(n\) is a finite run of expansion
steps from \(T^{(n)}\) to \(T^{(n+1)}\), and concatenating the stages gives an
expansion sequence from \(T^{(0)}=\langle B_\Gamma\rangle\).

\emph{Clause 2.}  Suppose the procedure never halts.  A closed branch admits
no step, so its only descendant is itself; hence the ancestor of an open
branch is open.  Let \(O_n\) be the set of open branches of \(T^{(n)}\),
nonempty because the procedure does not halt at stage \(n\).

First, infinitely many steps are performed.  Otherwise there is \(n_0\) such
that \(T^{(n)}=T^{(n_0)}\) for all \(n\geq n_0\).  Take \(B\in O_{n_0}\).  As
the procedure does not halt at stage \(n_0\) by its second test, \(B\) is not
saturated, so some task \(\tau\) available on \(B\) is not done.  If \(\tau\)
involves a term \(t_k\), let \(n\geq\max(n_0,k)\); otherwise let \(n=n_0\).
Then \(\tau\in W_n(B)\), and \(B\) is an open descendant of itself lacking
\(\tau\), so a step is performed at stage \(n\), a contradiction.

Next, the infinite branch construction.  Call \(B\in O_n\) \emph{good} if
for every \(n'\geq n\) some member of \(O_{n'}\) descends from \(B\).  The
unique member \(B_\Gamma\) of \(O_0\) is good, since every branch descends
from it.  Let \(B\in O_n\) be good, and let \(C_1,\ldots,C_q\) be the members
of \(O_{n+1}\) that descend from \(B\); there is at least one, because \(B\)
is good.  Every member of \(O_{n'}\), for \(n'\geq n+1\), that descends from
\(B\) descends from its own ancestor in \(T^{(n+1)}\), which is open and
descends from \(B\), so it is some \(C_p\).  If no \(C_p\) were good, each
\(C_p\) would have some \(n_p\) with no member of \(O_{n_p}\) descending from
it, and then none of any later \(O_{n'}\) either, since such a member would
have an open ancestor in \(T^{(n_p)}\) descending from \(C_p\).  For
\(n'\geq\max_p n_p\), no member of \(O_{n'}\) would descend from \(B\),
contradicting goodness.  So some \(C_p\) is good.  Let \(B^{0}=B_\Gamma\) and
let \(B^{n+1}\) be the first good member of \(O_{n+1}\) descending from
\(B^{n}\).

Let stage \(n\) begin at index \(j_n\) of the expansion sequence, so
\(T_{j_n}=T^{(n)}\); the \(j_n\) are nondecreasing and unbounded because
infinitely many steps occur.  For each index \(j\), choose \(n\) with
\(j_n\geq j\) and let \(B_j\) be the ancestor of \(B^{n}\) in \(T_j\).  It
does not depend on the choice of \(n\), because \(B^{n}\) is the ancestor of
\(B^{n+1}\) and ancestors compose.  Thus \(B_j\) is the parent of
\(B_{j+1}\), \((B_j)\) is a path, and every \(B_j\) is open.  By
\cref{thm:p02-w02-p02-014-fol-invariants} its limit \(B_\omega\) is open.

\(B_\omega\) is saturated.  Let \(\tau\) be available on \(B_\omega\).  Its
one or two positions lie on some \(B^{n_0}\).  Let \(n_1=\max(n_0,k)\) if
\(\tau\) involves \(t_k\), and \(n_1=n_0\) otherwise.  Since \(B^{n_1}\)
extends the list of \(B^{n_0}\), \(\tau\in W_{n_1}(B^{n_1})\).  When stage
\(n_1\) finishes processing \(\tau\), no open descendant of \(B^{n_1}\) lacks
\(\tau\); the ancestor of \(B^{n_1+1}\) at that moment is such a descendant
and is open, so \(\tau\) is in its done set, hence in that of
\(B^{n_1+1}\) and of \(B_\omega\).  By the remark after
\cref{def:p02-w02-p02-014-fol-branch}, a saturated branch is fair.

\(B_\omega\) is infinite.  Each step appends at least one sentence.  If
\(B_\omega\) were finite, the lists of the \(B^{n}\) would be constant from
some \(n_2\) on, so no step would replace a branch on the path after stage
\(n_2\) begins, and \(B_\omega=B^{n_2}\) as a pair of list and done set.  Then
\(B^{n_2}\) would be a finite, open and saturated branch of \(T^{(n_2)}\), and
the procedure would halt at stage \(n_2\), a contradiction.

\emph{Clause 3.}  If the procedure halts with \emph{unsatisfiable}, its output
\(T^{(n)}\) is the last member of a finite expansion sequence and is closed,
so \(\Gamma\) is unsatisfiable by \cref{thm:p02-w02-p02-014-fol-soundness}(3).
If it halts with \emph{satisfiable}, its output is an open saturated finite
branch, and \cref{thm:p02-w02-p02-014-fol-extraction} gives a model of
\(\Gamma\).  If it never halts, clause 2 and
\cref{thm:p02-w02-p02-014-fol-extraction} give a model of \(\Gamma\), namely
the expansion of \(\mathcal M_{\operatorname{fm}(B_\omega)}\restriction\mathcal
L_{\Gamma}\).  Hence for unsatisfiable \(\Gamma\) the only possibility is to
halt with \emph{unsatisfiable}, and conversely that output occurs only for
unsatisfiable \(\Gamma\).

\emph{Clause 4.}  The procedure is effective.  A computable enumeration of
\(\mathrm{CT}_{\Gamma}\) exists: for \(N=0,1,2,\ldots\), list the finitely
many closed terms of depth at most \(N\) built from the finitely many symbols
of \(\mathcal L_{\Gamma}\) and the parameters \(a_0,\ldots,a_N\), skipping
terms already listed; every closed term has some depth and uses finitely many
parameters, so it is listed.  The kind of a sentence is found by unique
readability (\cref{thm:fol-unique-readability}); substitution is defined by
recursion on complexity and is computable; the least unused parameter, the
occurrences of a term in an atomic sentence, and the closure test are finite
searches over a finite branch.  The saturation test for a finite branch is
decidable.  In \(\mathsf T^{=}\) no finite branch is saturated.  In
\(\mathsf T\) a finite branch carrying a gamma sentence is not saturated,
since its done set is finite while infinitely many tasks are available on
it; a finite branch carrying none is saturated exactly when
\(\langle i\rangle\) is done for every position \(i\) holding a
propositional-expandable or delta sentence, because these are then all its
available tasks.  With clause 3, the procedure halts with
\emph{unsatisfiable} on exactly the unsatisfiable finite sets.
\end{proof}
\end{annexitem}

\begin{annexitem}{annex:p02-w02-p02-014-fol-completeness}{thm:p02-w02-p02-014-fol-completeness}{Completeness and consequence}
\begin{proof}
\emph{Clause 1.}  A finite closed tableau rooted at \(\Gamma\) shows
\(\Gamma\) unsatisfiable by \cref{thm:p02-w02-p02-014-fol-soundness}(3).
Conversely, if \(\Gamma\) is unsatisfiable, the procedure of
\cref{def:p02-w02-p02-014-fol-procedure} halts with a finite closed tableau
rooted at \(\Gamma\) by \cref{thm:p02-w02-p02-014-fol-semidecision}(3).

\emph{Clause 2.}  By the definition of semantic consequence for sentences
(\cref{sec:fol-semantics-02}), \(\Gamma\entails\psi\) iff every model of
\(\Gamma\) satisfies \(\psi\).  If \(\mathcal M\entails\Gamma\cup\{\neg\psi\}\),
then \(\mathcal M\) is a model of \(\Gamma\) in which \(\psi\) is false, so
\(\Gamma\entails\psi\) fails.  If \(\Gamma\entails\psi\) fails, some model of
\(\Gamma\) falsifies \(\psi\) and therefore satisfies \(\neg\psi\).  So
\(\Gamma\entails\psi\) iff \(\Gamma\cup\{\neg\psi\}\) is unsatisfiable, iff
a finite closed tableau rooted at \(\Gamma\cup\{\neg\psi\}\) exists by
clause 1.  Every structure satisfies \(\varnothing\), so \(\entails\psi\) is
the case \(\Gamma=\varnothing\), with root \(\{\neg\psi\}\).

\emph{Clause 3.}  Apply \cref{thm:p02-w02-p02-014-fol-extraction} to the root
\(\Gamma\cup\{\neg\psi\}\): the structure obtained satisfies \(\Gamma\) and
\(\neg\psi\), so it falsifies \(\psi\).
\end{proof}
\end{annexitem}
